% Part: first-order-logic
% Chapter: introduction
% Section: soundness-completeness

\documentclass[../../../include/open-logic-section]{subfiles}

\begin{document}

\olfileid{fol}{int}{scp}

\olsection{Correctheid en volledigheid}

We voeren ook !!{derivation}ssystemen voor de eerste-ordelogica in.
Logici hebben vele !!{derivation}ssystemen ontwikkeld, maar zij
definiëren allemaal dezelfde !!{derivability}srelatie tussen
!!{sentence}s. We zeggen dat uit~$\Gamma$ de formule~$!A$ valt te
\emph{!!{derive}}, en schrijven $\Gamma \Proves !A$, als er een
!!a{derivation} van een bepaalde, precies gedefinieerde soort bestaat.
!!^{derivation}s zijn altijd eindige rangschikkingen van symbolen:
bijvoorbeeld een lijst !!{sentence}s of een ingewikkeldere structuur.
Het doel van !!{derivation}ssystemen is een middel te bieden om te
bepalen of !!a{sentence} een semantisch gevolg is van een
verzameling~$\Gamma$. Om dat doel te dienen moet $\Gamma \Entails !A$
dan en slechts dan gelden als $\Gamma \Proves !A$.

Als $\Gamma \Proves !A$ maar niet $\Gamma \Entails !A$, is ons
!!{derivation}ssysteem te sterk: het bewijst te veel. De eigenschap dat
uit $\Gamma \Proves !A$ steeds $\Gamma \Entails !A$ volgt, heet
\emph{correctheid}. Zij is een minimale eis voor ieder deugdelijk
!!{derivation}ssysteem. Als daarentegen $\Gamma \Entails !A$ maar niet
$\Gamma \Proves !A$, is ons !!{derivation}ssysteem te zwak: het bewijst
niet genoeg. De eigenschap dat uit $\Gamma \Entails !A$ steeds
$\Gamma \Proves !A$ volgt, heet \emph{volledigheid}. Correctheid is
doorgaans betrekkelijk eenvoudig te bewijzen, door inductie naar de
structuur van de inductief gedefinieerde !!{derivation}s. Volledigheid
is moeilijker te bewijzen.

Correctheid en volledigheid hebben enkele belangrijke gevolgen. Als uit
een verzameling !!{sentence}s~$\Gamma$ een tegenspraak, zoals
$!A \land \lnot !A$, valt te !!{derive}, heet zij \emph{inconsistent}.
Een inconsistente~$\Gamma$ kan geen model hebben en is dus
onvervulbaar. Uit volledigheid volgt het omgekeerde: iedere~$\Gamma$ die
niet inconsistent, oftewel \emph{consistent}, is, heeft een model. Dit
is zelfs equivalent met volledigheid en is de vorm van volledigheid die
we daadwerkelijk zullen bewijzen. Het is een diep en wellicht
verrassend resultaat: alleen al het feit dat men
$!A \land \lnot !A$ niet uit~$\Gamma$ kan bewijzen, garandeert dat er !!a{structure}
bestaat die is zoals~$\Gamma$ haar beschrijft. Volledigheid beantwoordt
daarmee de vraag welke verzamelingen !!{sentence}s modellen hebben:
precies de consistente verzamelingen.

De correctheids- en volledigheidsstelling hebben twee belangrijke
gevolgen: de compactheidsstelling en de stelling van
L\"owenheim--Skolem. Dit zijn belangrijke resultaten in de modeltheorie,
waaruit vele interessante resultaten kunnen worden afgeleid. Twee
ervan hebben we reeds genoemd: in de eerste-ordelogica kan niet worden
uitgedrukt dat het !!{domain} van !!a{structure} eindig of
!!{nonenumerable} is.

Historisch vergde dit alles veel tijd: de syntaxis en semantiek van de
eerste-ordelogica goed definiëren, deugdelijke !!{derivation}ssystemen
definiëren, hun correctheid en volledigheid bewijzen en duidelijk
krijgen wat eerste-ordetalen wel en niet kunnen uitdrukken. Inmiddels
weten we hoe dit moet, maar alle details doorwerken kan nog steeds
verwarrend en omslachtig zijn. Het is echter ook belangrijk, want de
hier ontwikkelde methoden voor de formele taal van de eerste-ordelogica
worden op grote schaal toegepast in de logica, informatica en
taalkunde. Het loont daarom op de lange duur de details door te werken.

\end{document}
